\documentclass[12pt,a4paper]{article}

% ---- Encoding & Font ----
\usepackage{iftex}
\ifPDFTeX
  \usepackage[utf8]{inputenc}
  \usepackage[T1]{fontenc}
\else
  \usepackage{fontspec}  % XeTeX/LuaTeX (tectonic): Unicode nativo (·, —, ¿, ª, §)
\fi
\usepackage[spanish]{babel}
\usepackage{csquotes}

% ---- Page layout ----
\usepackage[top=2.5cm, bottom=2.5cm, left=2.5cm, right=2.5cm]{geometry}
\usepackage{setspace}
\onehalfspacing

% ---- Math ----
\usepackage{amsmath, amssymb, amsthm}
\usepackage{mathtools}
\usepackage{physics}
\usepackage{esint}  % \oiint

% ---- Graphics ----
\usepackage{graphicx}
\usepackage{tikz}
\usepackage{float}
\usepackage{caption}

% ---- Tables ----
\usepackage{array}
\usepackage{booktabs}
\usepackage{tabularx}
\usepackage{colortbl}

% ---- Colours ----
\usepackage{xcolor}
\definecolor{notebg}{HTML}{FFF3CD}
\definecolor{noteborder}{HTML}{FFC107}
\definecolor{boxred}{HTML}{F8D7DA}
\definecolor{boxredborder}{HTML}{F5C6CB}

% ---- Boxes ----
\usepackage{mdframed}
\usepackage{tcolorbox}
\tcbuselibrary{most}

\newtcolorbox{keybox}[1][]{
  colback=blue!5,
  colframe=blue!70!black,
  arc=4pt,
  boxrule=1pt,
  fonttitle=\bfseries,
  title={#1}
}

\newtcolorbox{warnbox}[1][]{
  colback=boxred,
  colframe=boxredborder,
  coltitle=black!75,
  arc=4pt,
  boxrule=1pt,
  fonttitle=\bfseries,
  title={#1}
}

\newtcolorbox{notebox}[1][]{
  colback=notebg,
  colframe=noteborder,
  coltitle=black!75,
  arc=4pt,
  boxrule=1pt,
  fonttitle=\bfseries,
  title={#1}
}

% ---- Diagramas (gráficas y esquemas) ----
\usepackage{pgfplots}
\pgfplotsset{compat=1.18}
\usetikzlibrary{babel}  % babel español activa `>`; esto evita romper las flechas `->`
% courses/ElecMag2024/Clases/assets/tikzstyles.tex
% ---------------------------------------------------------------------------
% Estilos compartidos para los diagramas del curso Electromagnetismo (2024).
% Fuente única del "look" visual (colores, grosores, escalas) para que todas
% las figuras (TikZ / pgfplots / CircuiTikz) sean consistentes entre clases.
%
% Es una COPIA de courses/Fisica3-2015/Clases/assets/tikzstyles.tex: mismo
% dominio, misma paleta. Copia y no symlink para que cada curso sea
% autocontenido y la paleta de Física III pueda evolucionar aparte.
%
% Uso: la clase debe cargar antes `tikz` y `pgfplots` (y `circuitikz` si la
%      clase tiene circuitos), y luego % courses/ElecMag2024/Clases/assets/tikzstyles.tex
% ---------------------------------------------------------------------------
% Estilos compartidos para los diagramas del curso Electromagnetismo (2024).
% Fuente única del "look" visual (colores, grosores, escalas) para que todas
% las figuras (TikZ / pgfplots / CircuiTikz) sean consistentes entre clases.
%
% Es una COPIA de courses/Fisica3-2015/Clases/assets/tikzstyles.tex: mismo
% dominio, misma paleta. Copia y no symlink para que cada curso sea
% autocontenido y la paleta de Física III pueda evolucionar aparte.
%
% Uso: la clase debe cargar antes `tikz` y `pgfplots` (y `circuitikz` si la
%      clase tiene circuitos), y luego % courses/ElecMag2024/Clases/assets/tikzstyles.tex
% ---------------------------------------------------------------------------
% Estilos compartidos para los diagramas del curso Electromagnetismo (2024).
% Fuente única del "look" visual (colores, grosores, escalas) para que todas
% las figuras (TikZ / pgfplots / CircuiTikz) sean consistentes entre clases.
%
% Es una COPIA de courses/Fisica3-2015/Clases/assets/tikzstyles.tex: mismo
% dominio, misma paleta. Copia y no symlink para que cada curso sea
% autocontenido y la paleta de Física III pueda evolucionar aparte.
%
% Uso: la clase debe cargar antes `tikz` y `pgfplots` (y `circuitikz` si la
%      clase tiene circuitos), y luego \input{../assets/tikzstyles.tex}
% (compilar desde el directorio de la clase para que la ruta relativa resuelva).
% ---------------------------------------------------------------------------

% ---- Paleta (alineada con las cajas del preámbulo: keybox/notebox/warnbox) ----
\definecolor{figblue}{HTML}{1F4E79}   % azul principal (líneas, keybox)
\definecolor{figred}{HTML}{C0392B}    % rojo de énfasis (warnbox)
\definecolor{figamber}{HTML}{B8860B}  % ámbar (notebox)
\definecolor{figgray}{HTML}{555555}   % auxiliares, ejes, guías

% ---- CircuiTikz: escala y grosor de los componentes ----
% Guarda \ifdefined: la electrostática (Clases 1-14) no carga `circuitikz`, y
% sin la guarda el \input revienta con `Undefined control sequence \ctikzset`.
% Cuando el curso llegue a circuitos, el bloque se activa solo.
\ifdefined\ctikzset
  \ctikzset{
    bipoles/thickness=1,
    resistors/scale=0.9,
    capacitors/scale=0.9,
  }
\fi

% ---- TikZ: estilos reutilizables ----
% Nota: las puntas se declaran como `-{Latex[...]}` y no como `->`. La forma
% con llaves NO contiene un `>` literal, así que es inmune al `>` activo que
% introduce babel español (ver CLAUDE.md §7.3.3).
\usetikzlibrary{arrows.meta,calc,angles,quotes}

\tikzset{
  figline/.style={figblue, line width=0.9pt},
  figaux/.style={figgray, dashed, line width=0.6pt},
  % vectores
  figvec/.style={figblue, line width=0.9pt, -{Latex[length=2.2mm,width=1.6mm]}},
  figvecres/.style={figred, line width=1.3pt, -{Latex[length=2.6mm,width=1.9mm]}},
  figvecaux/.style={figgray, line width=0.7pt, -{Latex[length=1.8mm,width=1.3mm]}},
  figaxis/.style={figgray, line width=0.6pt, -{Latex[length=2mm,width=1.4mm]}},
  % cargas puntuales
  figcarga/.style={circle, inner sep=0pt, minimum size=5.4pt, draw=none},
  figpos/.style={figcarga, fill=figred},
  figneg/.style={figcarga, fill=figblue},
  figneutra/.style={figcarga, fill=figgray},
  % etiquetas
  figlbl/.style={font=\small},
  figlblaux/.style={font=\footnotesize, figgray},
}

% ---- pgfplots: estilo base de las gráficas ----
\pgfplotsset{
  emfig/.style={
    width=11cm, height=6.5cm,
    axis lines=left,
    xlabel style={font=\small},
    ylabel style={font=\small},
    tick label style={font=\footnotesize},
    every axis plot/.append style={line width=1pt},
    grid=major,
    grid style={figgray!20},
    legend cell align=left,
    legend style={font=\footnotesize, draw=figgray!40},
  },
}

% (compilar desde el directorio de la clase para que la ruta relativa resuelva).
% ---------------------------------------------------------------------------

% ---- Paleta (alineada con las cajas del preámbulo: keybox/notebox/warnbox) ----
\definecolor{figblue}{HTML}{1F4E79}   % azul principal (líneas, keybox)
\definecolor{figred}{HTML}{C0392B}    % rojo de énfasis (warnbox)
\definecolor{figamber}{HTML}{B8860B}  % ámbar (notebox)
\definecolor{figgray}{HTML}{555555}   % auxiliares, ejes, guías

% ---- CircuiTikz: escala y grosor de los componentes ----
% Guarda \ifdefined: la electrostática (Clases 1-14) no carga `circuitikz`, y
% sin la guarda el \input revienta con `Undefined control sequence \ctikzset`.
% Cuando el curso llegue a circuitos, el bloque se activa solo.
\ifdefined\ctikzset
  \ctikzset{
    bipoles/thickness=1,
    resistors/scale=0.9,
    capacitors/scale=0.9,
  }
\fi

% ---- TikZ: estilos reutilizables ----
% Nota: las puntas se declaran como `-{Latex[...]}` y no como `->`. La forma
% con llaves NO contiene un `>` literal, así que es inmune al `>` activo que
% introduce babel español (ver CLAUDE.md §7.3.3).
\usetikzlibrary{arrows.meta,calc,angles,quotes}

\tikzset{
  figline/.style={figblue, line width=0.9pt},
  figaux/.style={figgray, dashed, line width=0.6pt},
  % vectores
  figvec/.style={figblue, line width=0.9pt, -{Latex[length=2.2mm,width=1.6mm]}},
  figvecres/.style={figred, line width=1.3pt, -{Latex[length=2.6mm,width=1.9mm]}},
  figvecaux/.style={figgray, line width=0.7pt, -{Latex[length=1.8mm,width=1.3mm]}},
  figaxis/.style={figgray, line width=0.6pt, -{Latex[length=2mm,width=1.4mm]}},
  % cargas puntuales
  figcarga/.style={circle, inner sep=0pt, minimum size=5.4pt, draw=none},
  figpos/.style={figcarga, fill=figred},
  figneg/.style={figcarga, fill=figblue},
  figneutra/.style={figcarga, fill=figgray},
  % etiquetas
  figlbl/.style={font=\small},
  figlblaux/.style={font=\footnotesize, figgray},
}

% ---- pgfplots: estilo base de las gráficas ----
\pgfplotsset{
  emfig/.style={
    width=11cm, height=6.5cm,
    axis lines=left,
    xlabel style={font=\small},
    ylabel style={font=\small},
    tick label style={font=\footnotesize},
    every axis plot/.append style={line width=1pt},
    grid=major,
    grid style={figgray!20},
    legend cell align=left,
    legend style={font=\footnotesize, draw=figgray!40},
  },
}

% (compilar desde el directorio de la clase para que la ruta relativa resuelva).
% ---------------------------------------------------------------------------

% ---- Paleta (alineada con las cajas del preámbulo: keybox/notebox/warnbox) ----
\definecolor{figblue}{HTML}{1F4E79}   % azul principal (líneas, keybox)
\definecolor{figred}{HTML}{C0392B}    % rojo de énfasis (warnbox)
\definecolor{figamber}{HTML}{B8860B}  % ámbar (notebox)
\definecolor{figgray}{HTML}{555555}   % auxiliares, ejes, guías

% ---- CircuiTikz: escala y grosor de los componentes ----
% Guarda \ifdefined: la electrostática (Clases 1-14) no carga `circuitikz`, y
% sin la guarda el \input revienta con `Undefined control sequence \ctikzset`.
% Cuando el curso llegue a circuitos, el bloque se activa solo.
\ifdefined\ctikzset
  \ctikzset{
    bipoles/thickness=1,
    resistors/scale=0.9,
    capacitors/scale=0.9,
  }
\fi

% ---- TikZ: estilos reutilizables ----
% Nota: las puntas se declaran como `-{Latex[...]}` y no como `->`. La forma
% con llaves NO contiene un `>` literal, así que es inmune al `>` activo que
% introduce babel español (ver CLAUDE.md §7.3.3).
\usetikzlibrary{arrows.meta,calc,angles,quotes}

\tikzset{
  figline/.style={figblue, line width=0.9pt},
  figaux/.style={figgray, dashed, line width=0.6pt},
  % vectores
  figvec/.style={figblue, line width=0.9pt, -{Latex[length=2.2mm,width=1.6mm]}},
  figvecres/.style={figred, line width=1.3pt, -{Latex[length=2.6mm,width=1.9mm]}},
  figvecaux/.style={figgray, line width=0.7pt, -{Latex[length=1.8mm,width=1.3mm]}},
  figaxis/.style={figgray, line width=0.6pt, -{Latex[length=2mm,width=1.4mm]}},
  % cargas puntuales
  figcarga/.style={circle, inner sep=0pt, minimum size=5.4pt, draw=none},
  figpos/.style={figcarga, fill=figred},
  figneg/.style={figcarga, fill=figblue},
  figneutra/.style={figcarga, fill=figgray},
  % etiquetas
  figlbl/.style={font=\small},
  figlblaux/.style={font=\footnotesize, figgray},
}

% ---- pgfplots: estilo base de las gráficas ----
\pgfplotsset{
  emfig/.style={
    width=11cm, height=6.5cm,
    axis lines=left,
    xlabel style={font=\small},
    ylabel style={font=\small},
    tick label style={font=\footnotesize},
    every axis plot/.append style={line width=1pt},
    grid=major,
    grid style={figgray!20},
    legend cell align=left,
    legend style={font=\footnotesize, draw=figgray!40},
  },
}


% ---- Hyperlinks & TOC ----
\usepackage[hidelinks]{hyperref}
\usepackage{bookmark}

% ---- Enumerate ----
\usepackage{enumitem}

% ---- Miscellaneous ----
\usepackage{icomma}
\usepackage{siunitx}
\usepackage{textcomp}
\usepackage[bottom]{footmisc}

% ---- Title ----
\title{\textbf{Clase 6: La ecuación de Laplace en varias coordenadas} \\[0.3em]
        \large{Casos de una variable, simetría azimutal y separación de variables}}
\author{Transcripto y expandido --- Curso Electromagnetismo (OpenFING)}
\date{}

\begin{document}

\maketitle
\thispagestyle{empty}
\newpage

\tableofcontents
\newpage


\section{Dónde estamos}

Repaso de la clase anterior. La \textbf{ecuación de Poisson} corresponde al caso
en que hay una distribución continua de carga en un volumen; la \textbf{ecuación
de Laplace} es su ecuación homogénea asociada y vale cuando no hay densidad
volumétrica de carga ---toda la carga está en cargas puntuales o en superficies
que hacen de borde de la región---. De ella se probaron dos propiedades:

\begin{enumerate}[leftmargin=*]
  \item \textbf{Linealidad}: una combinación lineal de dos soluciones es
        solución, aunque \textbf{no tiene por qué cumplir las mismas condiciones
        de borde}.
  \item \textbf{Unicidad}: con condiciones de borde razonables (potencial dado, o
        componente normal de su gradiente dada), la solución es única a menos de
        una constante aditiva.
\end{enumerate}

\begin{notebox}[El plan de esta clase]
Todo eso era \textbf{formal}. El objetivo ahora es empezar a ser concreto y
producir soluciones: primero simplonas, después cada vez más sofisticadas. Pero
antes hace falta una herramienta previa.
\end{notebox}

\section{La ecuación de Laplace en otros sistemas de coordenadas}

\textbf{El motivo es práctico:} según la geometría del problema, las coordenadas
cartesianas no son siempre las más apropiadas.

\begin{figure}[H]
  \centering
  % Clase 6 — los dos sistemas de coordenadas, dibujados como en el pizarron.
% (a) cilindricas: r y theta se miden sobre la PROYECCION P' en el plano xy.
% (b) esfericas: r al punto, theta desde el eje z, phi azimutal sobre la proyeccion.
% Se dibujan juntos a proposito: es la unica forma de ver que e_r y theta no son
% lo mismo en un sistema y en el otro.
% Los versores van cortos y con el rotulo bien afuera: dibujados largos y con el
% rotulo pegado, se comen los nombres de P y P' (pasa en el render, no en el codigo).
\begin{tikzpicture}[scale=1.25]

% ---------------- (a) cilindricas ----------------
\begin{scope}
  \draw[figaxis] (0,0) -- (2.6,0);            \node[figlblaux] at (2.82,-0.02) {$x$};
  \draw[figaxis] (0,0) -- (-1.25,-1.05);      \node[figlblaux] at (-1.48,-1.22) {$y$};
  \draw[figaxis] (0,0) -- (0,2.5);            \node[figlblaux] at (-0.02,2.72) {$z$};
  \node[figlbl] at (-0.26,0.18) {$O$};
  \coordinate (Pp) at (1.5,-0.7);
  \coordinate (P)  at (1.5,1.15);
  \node[figneutra] at (Pp) {};
  \node[figneutra] at (P) {};
  \node[figlbl] at (1.28,1.38) {$P$};
  \node[figlbl] at (1.25,-0.95) {$P'$};
  \draw[figaux] (P) -- (Pp);
  % r sobre la proyeccion
  \draw[figvec] (0,0) -- (Pp);
  \node[figlbl, fill=white, inner sep=1.5pt] at (0.78,-0.2) {$r$};
  \draw[figgray, line width=0.6pt] (0.62,0) arc (0:-25:0.62);
  \node[figlblaux] at (0.9,-0.36) {$\theta$};
  % versores, cortos
  \draw[figvecres] (Pp) -- ++(0.5,-0.23);
  \node[figlbl, figred] at (2.34,-1.14) {$\hat e_r$};
  \draw[figvecres] (Pp) -- ++(0.23,0.5);
  \node[figlbl, figred] at (2.06,-0.02) {$\hat e_\theta$};
  \draw[figvecres] (P) -- ++(0,0.5);
  \node[figlbl, figred] at (1.5,1.88) {$\hat k$};
  \node[figlbl] at (0.6,-2.2) {(a) cil\'indricas};
  \node[figlblaux, align=center] at (0.6,-2.8)
    {$x=r\cos\theta$, $y=r\sin\theta$, $z=z$\\ $r$ y $\theta$ se miden sobre $P'$};
\end{scope}

% ---------------- (b) esfericas ----------------
\begin{scope}[xshift=6.9cm]
  \draw[figaxis] (0,0) -- (2.6,0);            \node[figlblaux] at (2.82,-0.02) {$x$};
  \draw[figaxis] (0,0) -- (-1.25,-1.05);      \node[figlblaux] at (-1.48,-1.22) {$y$};
  \draw[figaxis] (0,0) -- (0,2.5);            \node[figlblaux] at (-0.02,2.72) {$z$};
  \node[figlbl] at (-0.26,0.18) {$O$};
  \coordinate (Qp) at (1.5,-0.7);
  \coordinate (Q)  at (1.5,1.15);
  \node[figneutra] at (Qp) {};
  \node[figneutra] at (Qp) {};
  \node[figneutra] at (Q) {};
  \node[figlbl] at (1.25,1.4) {$P$};
  \node[figlbl] at (1.25,-0.95) {$P'$};
  \draw[figaux] (Q) -- (Qp);
  \draw[figaux] (0,0) -- (Qp);
  % r al punto, en el espacio
  \draw[figvec] (0,0) -- (Q);
  \node[figlbl, fill=white, inner sep=1.5pt] at (0.62,0.72) {$r$};
  \draw[figgray, line width=0.6pt] (0,0.95) arc (90:37.5:0.95);
  \node[figlblaux] at (0.36,1.12) {$\theta$};
  \draw[figgray, line width=0.6pt] (0.68,0) arc (0:-25:0.68);
  \node[figlblaux] at (0.96,-0.36) {$\varphi$};
  % versores, cortos y con el rotulo afuera
  \draw[figvecres] (Q) -- ++(0.4,0.31);
  \node[figlbl, figred] at (2.22,1.62) {$\hat e_r$};
  \draw[figvecres] (Q) -- ++(0.31,-0.4);
  \node[figlbl, figred] at (2.24,0.58) {$\hat e_\theta$};
  \draw[figvecres] (Q) -- ++(-0.2,-0.46);
  \node[figlbl, figred] at (0.78,0.5) {$\hat e_\varphi$};
  \node[figlbl] at (0.6,-2.2) {(b) esf\'ericas};
  \node[figlblaux, align=center] at (0.6,-2.8)
    {$z=r\cos\theta$, $x=r\sin\theta\cos\varphi$\\ $\theta$ se mide desde el eje $z$};
\end{scope}

\end{tikzpicture}

  \caption*{\small\itshape Los dos sistemas juntos, que es la única forma de ver
  que $\hat e_r$ y $\theta$ no significan lo mismo en uno y en otro.}
\end{figure}

\subsection{Cartesianas (repaso)}

\[
\grad\phi = \pdv{\phi}{x}\hat\imath + \pdv{\phi}{y}\hat\jmath + \pdv{\phi}{z}\hat k
\qquad
\div\vec V = \pdv{V_x}{x} + \pdv{V_y}{y} + \pdv{V_z}{z}
\]

\[
\laplacian\phi = \pdv[2]{\phi}{x} + \pdv[2]{\phi}{y} + \pdv[2]{\phi}{z}
\qquad\Longrightarrow\qquad
\text{Laplace: } \pdv[2]{\phi}{x} + \pdv[2]{\phi}{y} + \pdv[2]{\phi}{z} = 0
\]

(Poisson es lo mismo igualado a $-\rho/\varepsilon_0$.)

\subsection{Cilíndricas}

Se \textbf{mantiene $z$} como en cartesianas y en el plano $xy$ se pasa a
polares:

\[
x = r\cos\theta \qquad y = r\sin\theta \qquad z = z
\]

El punto de interés es $P$, en el espacio; lo que se usa para definir $r$ y
$\theta$ es su \textbf{proyección $P'$} sobre el plano $Oxy$. Los versores son
$\hat k$ (igual que antes), $\hat e_r$ ---radial, en la dirección de la
proyección horizontal--- y $\hat e_\theta$, ortogonal a él y en el sentido de los
$\theta$ crecientes.

\[
\grad\phi = \pdv{\phi}{r}\hat e_r + \frac{1}{r}\pdv{\phi}{\theta}\hat e_\theta
+ \pdv{\phi}{z}\hat k
\]

\[
\div\vec V = \frac{1}{r}\pdv{r}\big(r\,V_r\big) + \frac{1}{r}\pdv{V_\theta}{\theta}
+ \pdv{V_z}{z}
\]

Componiendo ambas:

\[
\boxed{\;\laplacian\phi = \frac{1}{r}\pdv{r}\left(r\pdv{\phi}{r}\right)
+ \frac{1}{r^2}\pdv[2]{\phi}{\theta} + \pdv[2]{\phi}{z}\;}
\]

\subsection{Esféricas}

Un punto $P$ se ubica con: la distancia al origen $r$; el ángulo $\theta$ que
forma con la \textbf{vertical} $z$; y el ángulo $\varphi$ que forma la proyección
$P'$ sobre el plano horizontal con el eje $x$.

\[
x = r\sin\theta\cos\varphi \qquad y = r\sin\theta\sin\varphi \qquad z = r\cos\theta
\]

\begin{notebox}[Sobre la letra $\varphi$]
Se usa para la coordenada azimutal, distinta del $\phi$ del potencial. El docente
reconoce la inconsistencia: es la notación más frecuente en la literatura, no la
que más le gusta.
\end{notebox}

Los tres versores son $\hat e_r$ (radial \textbf{en el espacio}, no en la
proyección), $\hat e_\theta$ (en el sentido de los $\theta$ crecientes) y
$\hat e_\varphi$ (en el de los $\varphi$ crecientes).

\[
\grad\phi = \pdv{\phi}{r}\hat e_r + \frac{1}{r}\pdv{\phi}{\theta}\hat e_\theta
+ \frac{1}{r\sin\theta}\pdv{\phi}{\varphi}\hat e_\varphi
\]

\[
\div\vec V = \frac{1}{r^2}\pdv{r}\big(r^2 V_r\big)
+ \frac{1}{r\sin\theta}\pdv{\theta}\big(\sin\theta\, V_\theta\big)
+ \frac{1}{r\sin\theta}\pdv{V_\varphi}{\varphi}
\]

\[
\boxed{\;\laplacian\phi = \frac{1}{r^2}\pdv{r}\left(r^2\pdv{\phi}{r}\right)
+ \frac{1}{r^2\sin\theta}\pdv{\theta}\left(\sin\theta\pdv{\phi}{\theta}\right)
+ \frac{1}{r^2\sin^2\theta}\pdv[2]{\phi}{\varphi}\;}
\]

\begin{warnbox}[Ojo con la notación]
$\hat e_r$ y $\hat e_\theta$ se llaman igual en cilíndricas y en esféricas pero
\textbf{significan cosas distintas}: el $\hat e_r$ cilíndrico apunta hacia la
proyección horizontal, el esférico apunta hacia el punto en el espacio; y el
$\theta$ cilíndrico se mide en el plano $xy$ mientras que el esférico se mide
desde el eje $z$. Es incómodo, pero es la costumbre de la literatura.
\end{warnbox}

\begin{notebox}[No hay que memorizarlas]
Hay una hoja de fórmulas de operadores diferenciales en el EVA que \textbf{se
puede llevar a los parciales}; el propio docente aclara que él tampoco se las
acuerda. Lo que sí hay que saber es que existen, y usarlas seguido.
\end{notebox}

\begin{notebox}[Un control barato: análisis dimensional]
Los factores $1/r$ tienen que estar porque $\theta$ y $\varphi$ son
\textbf{adimensionadas}, y una derivada respecto de ellas no tiene las mismas
unidades que una derivada respecto de $r$. El $\sin\theta$, en cambio, no sale
del análisis dimensional: hay que hacer la cuenta.
\end{notebox}

\subsection{Cómo se deducen: el gradiente en cilíndricas}

Ninguna de estas fórmulas cae del cielo: todas salen de aplicar la \textbf{regla
de la cadena} al cambio de variable. Se hace la más simple de todas como muestra;
las demás son idénticas conceptualmente, sólo que con muchas más cuentas.

Como $z$ no se toca, alcanza con ocuparse del plano $xy$. Del dibujo se leen las
relaciones entre versores:

\[
\hat e_r = \cos\theta\,\hat\imath + \sin\theta\,\hat\jmath
\qquad
\hat e_\theta = -\sin\theta\,\hat\imath + \cos\theta\,\hat\jmath
\]

\textbf{Paso 1: derivadas respecto de las coordenadas nuevas.} Por regla de la
cadena, y usando $x = r\cos\theta$, $y = r\sin\theta$:

\[
\pdv{\phi}{r} = \pdv{\phi}{x}\pdv{x}{r} + \pdv{\phi}{y}\pdv{y}{r}
= \cos\theta\,\pdv{\phi}{x} + \sin\theta\,\pdv{\phi}{y}
\]

\[
\pdv{\phi}{\theta} = \pdv{\phi}{x}\pdv{x}{\theta} + \pdv{\phi}{y}\pdv{y}{\theta}
= -r\sin\theta\,\pdv{\phi}{x} + r\cos\theta\,\pdv{\phi}{y}
\]

\textbf{Paso 2: invertir el sistema.} Es un sistema lineal de dos por dos en las
incógnitas $\partial_x\phi$ y $\partial_y\phi$ ---de hecho, dividiendo la segunda
por $r$, la matriz es una \textbf{rotación}---. Despejando:

\[
\pdv{\phi}{x} = \cos\theta\,\pdv{\phi}{r} - \frac{\sin\theta}{r}\pdv{\phi}{\theta}
\qquad
\pdv{\phi}{y} = \sin\theta\,\pdv{\phi}{r} + \frac{\cos\theta}{r}\pdv{\phi}{\theta}
\]

\begin{notebox}[Un error que apareció en el pizarrón y se corrigió sobre la marcha]
Estas dos expresiones estaban cruzadas ---la de $\partial_x$ escrita para
$\partial_y$ y viceversa---. El síntoma fue que al agrupar no aparecían los
versores esperados. Es el chequeo natural de esta cuenta: si al final no salen
$\hat e_r$ y $\hat e_\theta$, hay un error antes.
\end{notebox}

\textbf{Paso 3: sustituir y agrupar.} Partiendo del gradiente en cartesianas y
sustituyendo (el término en $\hat k$ sobrevive sin cambios):

\[
\grad\phi = \left(\cos\theta\,\partial_r\phi - \tfrac{\sin\theta}{r}\partial_\theta\phi\right)\hat\imath
+ \left(\sin\theta\,\partial_r\phi + \tfrac{\cos\theta}{r}\partial_\theta\phi\right)\hat\jmath
+ \partial_z\phi\,\hat k
\]

Agrupando por derivada en lugar de por versor, aparecen exactamente las
combinaciones del paso previo:

\begin{itemize}[nosep]
  \item el factor de $\partial_r\phi$ es
        $\cos\theta\,\hat\imath + \sin\theta\,\hat\jmath = \hat e_r$;
  \item el factor de $\tfrac{1}{r}\partial_\theta\phi$ es
        $-\sin\theta\,\hat\imath + \cos\theta\,\hat\jmath = \hat e_\theta$.
\end{itemize}

\[
\boxed{\;\grad\phi = \pdv{\phi}{r}\hat e_r + \frac{1}{r}\pdv{\phi}{\theta}\hat e_\theta
+ \pdv{\phi}{z}\hat k\;}
\]

\begin{notebox}[Y las demás igual]
Tomar las derivadas parciales del cambio de variable, sustituir y agrupar. «Tener
mucha paciencia, porque son muchas cuentas, pero al final no hay mayor misterio
que ése.»
\end{notebox}

\section{Laplace cuando hay una sola variable}

Se empieza por casos «bobos» ---matar una mosca con un cañón--- donde el potencial
depende de \textbf{una sola} coordenada. En todos ellos la EDP colapsa a una
ecuación diferencial \textbf{ordinaria}, y el resultado es algo que ya se conocía
de Física III.

\begin{notebox}[El valor del ejercicio es de verificación]
Comprobar que la ecuación de Laplace \textbf{da lo que tiene que dar} en las
situaciones ya conocidas. «Es como el primer \emph{check} en la lista de cosas a
verificar.»
\end{notebox}

\begin{figure}[H]
  \centering
  % Clase 6 — los tres casos de una sola variable y el perfil de potencial que da
% cada uno. El valor del ejercicio es de verificacion: Laplace reproduce lo que
% ya se sabia de Fisica III.
\begin{tikzpicture}

\begin{axis}[emfig, name=plano, width=4.15cm, height=4.2cm,
  xlabel={$x$}, ylabel={$\phi$}, domain=0:1, samples=2,
  xtick=\empty, ytick=\empty, ymin=0, ymax=1.15, xmin=0, xmax=1.1,
  title={\small cartesianas}, title style={font=\small}]
  \addplot[figblue] {0.15+0.85*x};
\end{axis}
\node[figlbl] at (2.05,-0.75) {$\phi = Ax+B$};
\node[figlblaux, align=center] at (2.05,-1.35) {condensador\\ de placas planas};

\begin{axis}[emfig, name=esf, at={(plano.east)}, anchor=west, xshift=0.8cm,
  width=4.15cm, height=4.2cm,
  xlabel={$r$}, ylabel={$\phi$}, domain=0.28:1.1, samples=80,
  xtick=\empty, ytick=\empty, ymin=0, ymax=3.9, xmin=0, xmax=1.2,
  title={\small esf\'ericas}, title style={font=\small}]
  \addplot[figblue] {1/x};
\end{axis}
\node[figlbl] at (7.0,-0.75) {$\phi = -A/r + B$};
\node[figlblaux, align=center] at (7.0,-1.35) {carga puntual,\\ esfera cargada};

\begin{axis}[emfig, name=cil, at={(esf.east)}, anchor=west, xshift=0.8cm,
  width=4.15cm, height=4.2cm,
  xlabel={$r$}, ylabel={$\phi$}, domain=0.12:1.1, samples=80,
  xtick=\empty, ytick=\empty, ymin=-2.6, ymax=0.6, xmin=0, xmax=1.2,
  title={\small cil\'indricas}, title style={font=\small}]
  \addplot[figblue] {ln(x)};
\end{axis}
\node[figlbl] at (11.9,-0.75) {$\phi = A\ln r + B$};
\node[figlblaux, align=center] at (11.9,-1.35) {l\'inea infinita,\\ condensador cil\'indrico};

\node[figlbl, align=center] at (7.0,-2.35)
  {En los tres casos la EDP colapsa a una EDO y las dos constantes\\[0.2em]
   quedan fijadas por los potenciales en los dos bordes.};

\end{tikzpicture}

  \caption*{\small\itshape Los tres perfiles conocidos, cada uno saliendo de la
  misma ecuación en el sistema de coordenadas que respeta la simetría.}
\end{figure}

\subsection{Cartesianas: \texorpdfstring{$\phi = \phi(x)$}{phi = phi(x)}}

Si $\phi$ no depende de $y$ ni de $z$, esas derivadas se anulan y queda

\[
\dv[2]{\phi}{x} = 0 \qquad\Longrightarrow\qquad \boxed{\;\phi(x) = A x + B\;}
\]

integrando dos veces. \textbf{Ejemplo físico: el condensador de placas
paralelas}, cuyas equipotenciales son planos. El potencial varía
\textbf{linealmente}, y los valores en las dos placas fijan $A$ y $B$: dos datos
para dos constantes.

\subsection{Esféricas: \texorpdfstring{$\phi = \phi(r)$}{phi = phi(r)}}

Si el problema tiene \textbf{simetría esférica} ---invariante bajo cualquier
rotación---, $\phi$ no depende de $\theta$ ni de $\varphi$ y esos dos términos del
laplaciano se anulan:

\[
\frac{1}{r^2}\dv{r}\left(r^2\dv{\phi}{r}\right) = 0
\]

El factor $1/r^2$ no juega ningún papel, así que $r^2\,\phi'(r) = A$, es decir
$\phi'(r) = A/r^2$, y primitivando:

\[
\boxed{\;\phi(r) = -\frac{A}{r} + B\;}
\]

\begin{notebox}[Ejemplos]
Una carga puntual, o una esfera uniformemente cargada vista desde afuera. El
resultado es el $1/r$ conocido, más la constante aditiva que siempre se puede
sumar. Imponiendo condiciones de borde se identifica $A$ con
$-Q/4\pi\varepsilon_0$ (a menos de un signo, según la convención con que se
escriba).
\end{notebox}

\subsection{Cilíndricas: \texorpdfstring{$\phi = \phi(r)$}{phi = phi(r)}}

Con \textbf{simetría cilíndrica} ---sin dependencia en $z$ ni en $\theta$---
sobrevive un solo término:

\[
\frac{1}{r}\dv{r}\left(r\dv{\phi}{r}\right) = 0
\qquad\Longrightarrow\qquad
r\,\phi'(r) = A
\qquad\Longrightarrow\qquad
\boxed{\;\phi(r) = A\ln r + B\;}
\]

\begin{warnbox}[Las hipótesis son restrictivas y hay que decirlas]
Los ejemplos son una \textbf{línea infinita de carga} (cuyo potencial logarítmico
ya se conocía) o un \textbf{condensador cilíndrico}. Pero para que no haya
dependencia en $z$ el sistema tiene que ser \textbf{muy largo}, y para que no
dependa de $\theta$ tiene que ser isótropo en el plano.
\end{warnbox}

\section{Problemas con simetría azimutal}

Ahora se sube un escalón: un problema que depende de \textbf{dos} variables. (El
caso de tres coordenadas no se ve en el teórico de este curso.)

A la variable $\varphi$ se la llama \textbf{azimutal}, y \textbf{simetría
azimutal} significa que el problema no depende de ella: es invariante bajo
rotaciones alrededor del eje $z$, aunque no bajo rotaciones arbitrarias.

\begin{figure}[H]
  \centering
  % Clase 6 — que significa simetria azimutal, y por que el problema de la esfera
% en campo uniforme la tiene. Lo que importa aca no es la polarizacion (eso es
% la figura de la Clase 5) sino cual rotacion deja el problema invariante.
% Los arcos de rotacion van DEBAJO de donde terminan las lineas de campo: puestos
% a su altura se cruzan con ellas y con sus propios rotulos.
\begin{tikzpicture}[scale=1]

% ---------------- (a) rotacion alrededor de z: invariante ----------------
\begin{scope}
  \draw[figaxis] (0,-1.35) -- (0,2.3);
  \node[figlblaux] at (0,2.52) {$z$};
  \foreach \xx in {-1.75,-1.1,1.1,1.75}{
    \draw[figvecaux] (\xx,-1.2) -- (\xx,1.85);
  }
  \node[figlblaux] at (-1.75,2.1) {$\vec E_0$};
  \fill[figgray!25] (0,0) circle (0.72);
  \draw[figline] (0,0) circle (0.72);
  \coordinate (P) at (1.15,1.28);
  \node[figneutra] at (P) {};
  \draw[figvec] (0,0) -- (P);
  \node[figlbl, fill=white, inner sep=1.5pt] at (0.82,0.98) {$r$};
  \draw[figgray, line width=0.6pt] (0,0.95) arc (90:48:0.95);
  \node[figlblaux] at (0.36,1.08) {$\theta$};
  % rotacion alrededor de z, debajo de todo
  \draw[figred, line width=0.9pt] (-0.9,-1.85) arc (180:360:0.9 and 0.3);
  \draw[figred, line width=0.9pt, -{Latex[length=2mm,width=1.5mm]}]
    (0.9,-1.85) arc (0:22:0.9 and 0.3);
  \node[figlbl, figred] at (0,-2.42) {rotar en $\varphi$};
  \node[figlbl] at (0,-3.05) {(a) invariante};
\end{scope}

% ---------------- (b) rotacion alrededor de otro eje: NO ----------------
\begin{scope}[xshift=6.9cm]
  \draw[figaxis] (0,-1.35) -- (0,2.3);
  \node[figlblaux] at (0,2.52) {$z$};
  \foreach \s in {-1,1}{
    \draw[figvecaux] ({\s*1.75-0.6},-1.2) -- ({\s*1.75+0.6},1.85);
    \draw[figvecaux] ({\s*1.1-0.6},-1.2) -- ({\s*1.1+0.6},1.85);
  }
  \node[figlblaux] at (-1.0,2.1) {$\vec E_0$ rotado};
  \fill[figgray!25] (0,0) circle (0.72);
  \draw[figline] (0,0) circle (0.72);
  % rotacion alrededor de un eje horizontal, tambien abajo
  \draw[figred, line width=0.9pt] (-0.9,-1.85) arc (180:360:0.9 and 0.3);
  \draw[figred, line width=0.9pt, -{Latex[length=2mm,width=1.5mm]}]
    (0.9,-1.85) arc (0:22:0.9 and 0.3);
  \node[figlbl, figred] at (0,-2.42) {rotar en otro eje};
  \node[figlbl] at (0,-3.05) {(b) el campo se mueve};
\end{scope}

\node[figlbl, align=center] at (3.45,-3.95)
  {La esfera es sim\'etrica bajo cualquier rotaci\'on; el campo uniforme, s\'olo
   bajo las de eje $z$.\\[0.2em]
   Por eso $\phi=\phi(r,\theta)$: depende de dos variables, no de una ni de tres.};

\end{tikzpicture}

  \caption*{\small\itshape La simetría no es del cuerpo ni del campo por
  separado, sino del problema completo: es la que sobrevive a los dos a la vez.}
\end{figure}

Con el último término del laplaciano anulado, la ecuación es

\[
\frac{1}{r^2}\pdv{r}\left(r^2\pdv{\phi}{r}\right)
+ \frac{1}{r^2\sin\theta}\pdv{\theta}\left(\sin\theta\pdv{\phi}{\theta}\right) = 0
\]

Siguen siendo derivadas parciales, porque siguen siendo dos variables.

\subsection{Buscar una base, no la solución general}

Acá entra la propiedad de linealidad de la clase anterior. Si el conjunto de
soluciones es un \textbf{espacio vectorial}, lo cómodo ---como en álgebra
lineal--- es describir sus elementos desarrollados en una \textbf{base}. Entonces
el objetivo no es hallar directamente la solución del problema concreto sino

\[
\phi(r,\theta) = \sum_n a_n\,\phi_n(r,\theta)
\]

donde las $\phi_n$ son una base de soluciones.

\begin{warnbox}[Esta suma no es la de álgebra lineal]
Allá las bases son finitas y cada elemento se escribe como suma finita. Acá el
espacio tiene \textbf{dimensión infinita} y esto es una \textbf{serie}, con un
proceso de convergencia detrás. Es una \textbf{base de Hilbert}, no una base en el
sentido usual.

El docente es explícito sobre cómo se trata eso en la práctica: «ustedes son
candidatos a ingenieros y yo soy físico, así que ese problema se lo dejamos a los
matemáticos», pero \textbf{hay que saber que el peligro existe}. Y agrega el
criterio pragmático: una base es útil si sumando unos diez términos ya se parece;
si hacen falta 800\,000, no sirve para nada. El objetivo es reemplazar un problema
complicado por uno manejable.
\end{warnbox}

\subsection{El método de separación de variables}

El truco para encontrar elementos de la base es buscar soluciones \textbf{con
variables separables}, es decir, de la forma producto:

\[
\phi(r,\theta) = Z(r)\,P(\theta)
\]

\begin{warnbox}[La advertencia central de la clase]
Las soluciones \textbf{en general no son de esta forma} ---típicamente las dos
variables aparecen entreveradas---. Lo que se afirma es mucho más débil: que se
puede encontrar \textbf{una base} formada por soluciones de este tipo. La solución
general es una \textbf{combinación lineal} de ellas, y no tiene forma de producto.
\end{warnbox}

\begin{notebox}[El método es mucho más general que este problema]
Sirve para \textbf{cualquier ecuación diferencial lineal homogénea}, y con más
trabajo incluso para algunas no lineales. La metáfora del docente: es una 4$\times$4
--- no cruza el océano, pero cruza muchísimos ríos. Vale la pena tenerlo en la caja
de herramientas.

(No confundir con el método de separación de variables de Cálculo 1. Se llama
igual y la cuenta se parece, pero es otra cosa.)
\end{notebox}

Al sustituir el producto, cada derivada golpea a una sola función y la otra sale
de factor:

\[
\frac{P(\theta)}{r^2}\dv{r}\big(r^2 Z'(r)\big)
+ \frac{Z(r)}{r^2\sin\theta}\dv{\theta}\big(\sin\theta\,P'(\theta)\big) = 0
\]

Nótese que ya son derivadas \textbf{ordinarias}: $Z$ depende sólo de $r$ y $P$
sólo de $\theta$.

\subsection{La constante de separación}

Hasta acá sólo se factorizó. La separación viene ahora: se multiplica por $r^2$ y
se divide por $Z\,P$, y se pasa un término al otro miembro:

\[
\frac{1}{Z(r)}\dv{r}\big(r^2 Z'(r)\big)
= -\frac{1}{P(\theta)\sin\theta}\dv{\theta}\big(\sin\theta\,P'(\theta)\big)
\]

\textbf{El argumento decisivo, planteado como pregunta capciosa: ¿de qué variable
depende cada miembro?} De ninguna.

\begin{itemize}[leftmargin=*]
  \item El miembro izquierdo tiene pinta de depender de $r$, pero es igual al
        derecho, que no depende de $r$.
  \item El derecho tiene pinta de depender de $\theta$, pero es igual al
        izquierdo, que no depende de $\theta$.
\end{itemize}

Luego ambos son iguales a una misma \textbf{constante}, la \textbf{constante de
separación} $K$:

\[
\boxed{\;\dv{r}\big(r^2 Z'(r)\big) = K\,Z(r)\;}
\qquad\qquad
\boxed{\;\dv{\theta}\big(\sin\theta\,P'(\theta)\big) + K\sin\theta\,P(\theta) = 0\;}
\]

\begin{keybox}[Lo que se ganó]
El problema de hallar una base de soluciones de una \textbf{ecuación en derivadas
parciales} quedó reducido a resolver \textbf{dos ecuaciones diferenciales
ordinarias}. Las EDO, que uno creía ya dominadas, pasan a ser el juego lateral
para resolver el problema grande.
\end{keybox}

Ambas ecuaciones son de \textbf{segundo orden}, así que para cada valor aceptable
de $K$ hay típicamente dos soluciones de cada una: cuatro combinaciones. Y a
priori hay una para cada $K$, es decir infinitas --- \textbf{aunque se verá que no
todos los valores de $K$ son aceptables}, y por eso al final no serán tantas.

\section{La ecuación de Legendre}

Se ataca primero la \textbf{ecuación angular}.

\subsection{El cambio de variable \texorpdfstring{$u = \cos\theta$}{u = cos theta}}

Como toda la ecuación es trigonométrica, conviene pasar a una variable
trigonométrica:

\[
u = \cos\theta \qquad\Longrightarrow\qquad \dv{u}{\theta} = -\sin\theta
\]

Por regla de la cadena,
$\dv{P}{\theta} = \dv{P}{u}\dv{u}{\theta} = -\sin\theta\,\dv{P}{u}$, de modo que
la combinación que aparece dentro de la derivada es

\[
\sin\theta\,\dv{P}{\theta} = -\sin^2\theta\,\dv{P}{u} = (u^2 - 1)\dv{P}{u}
\]

usando $\sin^2\theta = 1 - \cos^2\theta = 1 - u^2$. Derivando otra vez respecto de
$\theta$ ---lo que vuelve a traer un factor $-\sin\theta$ de la cadena---:

\[
\dv{\theta}\left(\sin\theta\dv{P}{\theta}\right)
= -\sin\theta\,\dv{u}\left[(u^2-1)\dv{P}{u}\right]
\]

Sustituyendo en la ecuación angular, \textbf{el $\sin\theta$ se cancela} y queda
una ecuación mucho más limpia:

\[
\boxed{\;\dv{u}\left[(1-u^2)\dv{P}{u}\right] + K\,P(u) = 0\;}
\]

que es la \textbf{ecuación de Legendre}. Distribuyendo la derivada se obtiene la
forma equivalente

\[
(1-u^2)\,P''(u) - 2u\,P'(u) + K\,P(u) = 0
\]

\begin{warnbox}[Por qué esto no es fácil]
Es una EDO lineal de segundo orden, sí, \textbf{pero no de coeficientes
constantes}: los coeficientes dependen de $u$. Con coeficientes constantes ---como
la ecuación del resorte--- las soluciones serían exponenciales o senos y cosenos.
Acá no: resolverla da trabajo real.
\end{warnbox}

\subsection{Qué soluciones interesan}

Antes de resolverla hay que acotar qué se busca:

\begin{itemize}[leftmargin=*]
  \item \textbf{Soluciones suaves, sin singularidades.} Es coherente con el
        planteo: se está aplicando la ecuación de Laplace, es decir se está en una
        zona \textbf{sin cargas}, así que no debe haber singularidades.
  \item \textbf{$\theta$ recorriendo todo el rango $[0,\pi]$}, extremos incluidos.
        Equivalentemente, dado que $\cos 0 = 1$ y $\cos\pi = -1$:
        \[
        \theta \in [0,\pi] \qquad\Longleftrightarrow\qquad u \in [-1,1]
        \]
\end{itemize}

\begin{figure}[H]
  \centering
  % Clase 6 — que problemas quedan dentro del analisis y cuales no. El casquete
% esferico tiene simetria azimutal igual, pero theta no recorre todo [0,pi], y
% eso cambia el conjunto de soluciones de Legendre.
% El barrido de theta va por FUERA de la esfera y punteado: dibujado encima se
% confunde con el radio y con el propio contorno.
\begin{tikzpicture}[scale=1]

% ---------------- (a) esfera completa: theta en [0,pi] ----------------
\begin{scope}
  \draw[figaxis] (0,-1.75) -- (0,2.1); \node[figlblaux] at (0,2.32) {$z$};
  \draw[figline] (0,0) circle (1.15);
  \draw[figaux] (0,0) ellipse (1.15 and 0.37);
  \node[figneutra] at (0,0) {};
  \draw[figvecres] (0,0) -- (58:1.15);
  \draw[figred, line width=0.8pt] (0,0.62) arc (90:58:0.62);
  \node[figlblaux, figred] at (0.3,0.76) {$\theta$};
  % barrido, por fuera del contorno
  \draw[figgray, dashed, line width=0.8pt, -{Latex[length=2mm,width=1.5mm]}]
    (0,1.62) arc (90:-90:0.42 and 1.62);
  \node[figlblaux, fill=white, inner sep=1.5pt] at (1.15,-1.42) {$\theta$ barre};
  \node[figlbl] at (0,-2.3) {(a) $\theta\in[0,\pi]$};
  \node[figlbl, figblue] at (0,-2.9) {$u=\cos\theta\in[-1,1]$};
\end{scope}

% ---------------- (b) casquete: theta en [0,alfa] ----------------
\begin{scope}[xshift=6.6cm]
  \draw[figaxis] (0,-1.75) -- (0,2.1); \node[figlblaux] at (0,2.32) {$z$};
  \node[figneutra] at (0,0) {};
  % casquete: arco entre 40 y 140 grados
  \draw[figline, line width=1.3pt] (40:1.15) arc (40:140:1.15);
  \draw[figaux] (0,0) -- (40:1.15);
  \draw[figaux] (0,0) -- (140:1.15);
  \draw[figaux] (0,0.881) ellipse (0.739 and 0.24);
  \draw[figred, line width=0.8pt] (0,0.62) arc (90:40:0.62);
  \node[figlblaux, figred, fill=white, inner sep=1.2pt] at (0.3,0.42) {$\alpha$};
  % barrido acotado, tambien por fuera
  \draw[figgray, dashed, line width=0.8pt, -{Latex[length=2mm,width=1.5mm]}]
    (0,1.62) arc (90:40:0.42 and 1.62);
  \node[figlblaux, fill=white, inner sep=1.5pt] at (1.78,0.42) {s\'olo hasta $\alpha$};
  \draw[figvecaux] (1.32,0.5) -- (0.72,0.92);
  \node[figlbl] at (0,-2.3) {(b) $\theta\in[0,\alpha]$};
  \node[figlbl, figred] at (0,-2.9) {queda fuera del an\'alisis};
\end{scope}

% Pie en tres lineas cortas: el texto de los nodos NO escala con `scale`, asi que
% una linea larga desborda el margen aunque se achique la geometria (CLAUDE.md 6.5.6).
\node[figlbl, align=center] at (3.3,-3.75)
  {Los dos tienen simetr\'ia azimutal. Pero cuanto m\'as\\[0.2em]
   restringido est\'a el recorrido de $u$, \emph{m\'as} soluciones\\[0.2em]
   admite Legendre: pedir el intervalo completo es lo\\[0.2em]
   que deja el problema simple.};

\end{tikzpicture}

  \caption*{\small\itshape La restricción no es sobre la simetría sino sobre el
  recorrido de $\theta$, y es lo que decide cuántas soluciones admite Legendre.}
\end{figure}

\begin{notebox}[Qué queda afuera, y por qué se dice]
Un problema con simetría azimutal donde $\theta$ sólo recorre $[0,\alpha]$ ---por
ejemplo un \textbf{casquete esférico}, un cascarón limitado por un cono de
semiángulo $\alpha$--- \textbf{no está incluido} en este análisis. No es que no se
pueda: se resuelve con los mismos métodos, pero \textbf{hay más soluciones
posibles} y hay que trabajar más.

La regla general es ésa: \textbf{cuanto más restringido está el recorrido de $u$,
más soluciones hay}. Pedir el intervalo completo $[-1,1]$ es la exigencia más
fuerte y por eso deja el conjunto de soluciones más chico --- y el problema más
simple.
\end{notebox}

\bigskip
\noindent\emph{La clase termina sin resolver la ecuación de Legendre; queda para
la Clase 7, junto con la ecuación radial.}

\end{document}
